% Generated draft: requires mathematical and source review.
\documentclass[11pt]{article}
\usepackage[margin=25mm]{geometry}
\usepackage{fontspec}
\setmainfont{Latin Modern Roman}
\usepackage{amsmath,amssymb,mathtools,unicode-math}
\setmathfont{Latin Modern Math}
\usepackage{enumitem,needspace,xurl,hyperref}
\setlength{\parindent}{0pt}
\setlength{\parskip}{5pt}
\setlength{\emergencystretch}{3em}
\newcommand{\N}{\mathbb{N}}
\newcommand{\Z}{\mathbb{Z}}
\newcommand{\Q}{\mathbb{Q}}
\newcommand{\R}{\mathbb{R}}
\newcommand{\C}{\mathbb{C}}
\newcommand{\F}{\mathbb{F}}
\newcommand{\A}{\mathbb{A}}
\newcommand{\PP}{\mathbb{P}}
\newcommand{\E}{\mathbb{E}}
\newcommand{\eps}{\varepsilon}
\newcommand{\dd}{\,\mathrm{d}}
\begin{document}
\section*{UnsolvedMath ID: 701025}
\textbf{Source problem code:} catalogue entry
\subsection{Definitions and mathematical statement}
Let \(p\) be a prime. A pro-\(p\) group is a topological group isomorphic to an inverse limit of finite \(p\)-groups. A free pro-\(p\) group \(F\) has a topological generating basis \(X\) converging to the identity, with the following universal property: for every pro-\(p\) group \(H\), each map \(f:X\to H\) converging to the identity extends uniquely to a continuous homomorphism \(F\to H\). Here convergence means that every neighbourhood of the identity contains all but finitely many of the specified elements; this condition is automatic for finite \(X\). The rank is \(|X|\).

Write \(F_{\mathrm{abs}}\) for the underlying abstract group of \(F\). Interpreting representability by matrices as a faithful representation by matrices of some finite size, the question is whether there exist an integer \(n\geq 1\), a commutative associative ring \(R\) with identity, and an injective group homomorphism
\[
\rho:F_{\mathrm{abs}}\hookrightarrow \operatorname{GL}_n(R),
\]
where \(\operatorname{GL}_n(R)\) is the group of invertible \(n\times n\) matrices over \(R\). No continuity condition is imposed on \(\rho\), and the source imposes no characteristic restriction on \(R\).

The source does not specify the rank of \(F\), whether infinite rank is included, or whether the question is intended uniformly for every prime and every rank. Thus the displayed existence question records the matrix-representation target, while its intended range of quantification remains unresolved.
\subsection{Short English statement}
Can the underlying abstract group of a free pro-\(p\) group be embedded into a group of finite square matrices over a commutative associative ring with identity? The source leaves the prime and rank scope unspecified.
\subsection{Sources}
\begin{itemize}
\item Local UnsolvedMath registry: supplied complete source record, permanent ID \(701025\), exact source problem code \texttt{catalogue entry}, title \emph{Kourovka Notebook Problem 14.40}; upstream revision \texttt{d371ee917ac9647fbb06cd6fc529deec6f3dda82}.
\item A. N. Zubkov, Problem 14.40, \emph{The Kourovka Notebook (Unsolved Problems in Group Theory)}, 14th issue (1999), as attributed in the supplied record. The background identifies the 21st issue edition (2026), arXiv:1401.0300v48. The supplied collection description names E. I. Khukhro and V. D. Mazurov as editors.
\item The registry labels the problem \emph{open}; this is supplied source metadata, not an independently verified mathematical or literature assessment. The record gives an update timestamp of \texttt{2026-10-16T14:00:00Z}, but supplies no literature-verification date or literature-source links. The cited edition has not been independently consulted for this draft.
\end{itemize}
\end{document}
